% II. Анализ на възможните решения и обосновка на избора.
% Правило за раздела: всяка алтернатива се описва с това, което тя прави ДОБРЕ, и чак
% след това с причината да бъде отхвърлена. Отхвърляне без изложено предимство е
% реклама на избраното решение, не анализ.
\section{Анализ на възможните решения и обосновка на избора}
\label{sec:alternatives}

Проектът има четири независими проектни решения: език и среда за изпълнение, модел на
паралелизъм, начин на представяне на задачата пред алгоритъма и стратегия за обмен на
информация между нишките. Разделът разглежда алтернативите по всяко от тях и
обосновава направения избор. Обосновката е предварителна там, където зависи от
измерване, и в такива случаи е отбелязана изрично.

\subsection{Език и среда за изпълнение}
\label{subsec:alt-language}

Разгледани са три възможности. Език с управлявана среда за изпълнение би съкратил
времето за разработка, но въвежда събиране на боклука в най-горещия цикъл на
търсенето и прави измерването на времето за една итерация зависимо от фактор извън
контрола на програмата. Език с явно управление на паметта и без среда за изпълнение
дава пълен контрол, но за част от нужните конструкции няма зряла стандартна
поддръжка. Избран е C\texttt{++} по стандарта ISO/IEC 14882:2020~\cite{iso_cpp20},
защото дава едновременно контрол върху разположението на данните в паметта и
стандартни средства за паралелизъм, тоест позволява измерването да отразява
алгоритъма, а не средата под него.

\subsection{Модел на паралелизъм}
\label{subsec:alt-parallelism}

Първата възможност е директно управление на нишки. Тя дава най-точен контрол върху
свързването на нишка с ядро и върху момента на синхронизация, но прехвърля цялата
координация върху програмиста. Втората е задачно ориентирана среда с крадене на
работа, която се справя по-добре с неравномерно натоварване, но добавя слой, чиято
цена трябва да бъде отчетена отделно при измерването. Третата е директивен подход над
цикли, който е най-икономичен откъм код, но се вписва зле в търсене с нерегулярна
структура.

Избран е първият вариант като основен, с явно управление на нишките и с атомарни
операции за споделеното състояние, защото предметът на измерването е именно цената на
синхронизацията и тя не бива да бъде скрита зад чужд планировчик.
\TODO{Да се реши дали задачно ориентираната среда влиза като втори, сравнителен
вариант, или остава в предложенията за по-нататъшна работа.}

\subsection{Представяне на задачата пред алгоритъма}
\label{subsec:alt-problem}

Възможните решения са две. Динамичен полиморфизъм чрез виртуални функции дава един
общ договор и позволява алгоритъмът да бъде компилиран веднъж, но вкарва непряко
извикване в най-често изпълнявания път, тоест в оценяването на целевата функция.
Статичен полиморфизъм чрез шаблони и концепции премахва това извикване и отваря
целевата функция за вграждане и за векторизация, но увеличава времето за компилация и
размера на кода.

Избран е статичният вариант за горещия път и общ динамичен договор на нивото на
конфигуриране на експеримента. Разделянето на двете нива е проектно решение, а не
компромис: изборът на алгоритъм се прави веднъж на изпълнение, а оценяването на
целевата функция се прави милиони пъти.

\subsection{Стратегия за обмен между нишките}
\label{subsec:alt-cooperation}

Трите разгледани схеми следват класификацията от~\cite{crainic2010parallel}:
независим многократен старт, кооперативна схема със споделено най-добро решение и
островен модел с периодична миграция. Те не са взаимно изключващи се, а образуват
редица с нарастваща комуникация, затова и трите са реализирани и влизат в
сравнението. Твърдението, че по-голямата комуникация се изплаща с по-добро качество,
е хипотеза за проверка, а не изходно допускане.

\subsection{Обобщение на избора}
\label{subsec:alt-summary}

\TODO{Таблица с колони: проектно решение, разгледани алтернативи, избрано решение,
основание. Попълва се, след като изборът по §2.2 бъде окончателно потвърден.}
